% !TEX program = xelatex
% 编译方式：在本目录执行两遍 xelatex（详见 docs/_manual_common/README.md）。
% 源码依据：src/graph/ 与 include/hacdcpf/graph/
%   （topology_analysis / switch_contraction / series_reduction / kron_reduction / result_recovery），公式与参数以代码为准。
\documentclass[UTF8,fontset=fandol,a4paper,11pt]{ctexart}
\usepackage{../../_manual_common/hysim_manual}

\title{\textbf{交直流混合高弹性能源电力系统仿真平台}\\[0.5em]
       {\Large 图建模、拓扑分析与网络降阶技术手册}\\[0.3em]
       {\large 数学模型、约束、算法流程、数值稳定性与工业级评估}}
\author{HySim-XJTU-HRPES（西安交通大学 HRPES 团队）}
\date{\today}

\begin{document}
\maketitle
\begin{abstract}
\noindent 本手册依据 HySim-XJTU-HRPES 平台图模块
（\srcpath{src/graph/}、\srcpath{include/hacdcpf/graph/}）整理，覆盖从
\texttt{HybridPowerSystem} 到 \texttt{PowerSystemGraph} 的图建模、拓扑分析（孤岛/辐射性/
回路/桥/割点），以及零阻抗合并与开关收缩、串联/悬垂降阶、稠密与稀疏 Kron/Schur 补降阶，
并为每类降阶保留映射以恢复被消去母线的状态。全部结论以源码为准：Kron 仅对无源/线性内部
节点精确，悬垂折叠为近似；AC/DC 母线共享整数 ID 时须使用 domain-qualified 映射
（\fld{ac\_bus\_id\_to\_node\_idx}/\fld{dc\_bus\_id\_to\_node\_idx}）。伞形头文件为
\srcpath{hacdcpf/graph/graph.hpp}。文档风格、参数表体例与符号约定与《潮流计算技术手册》
《最优潮流技术手册》一致。
\end{abstract}

\tableofcontents
\newpage

\section{阅读指南与约定}
\subsection{数据来源}
\begin{itemize}
  \item \srcpath{include/hacdcpf/graph/power\_system\_graph.hpp} —— 图建模与
        \texttt{PowerSystemGraph}/\allowbreak\texttt{GraphNode}/\allowbreak\texttt{GraphEdge}；
  \item \srcpath{include/hacdcpf/graph/topology\_analysis.hpp} —— 拓扑分析与
        \texttt{TopologyReport}；
  \item \srcpath{include/hacdcpf/graph/switch\_contraction.hpp}、
        \srcpath{series\_reduction.hpp}、\srcpath{kron\_reduction.hpp}、
        \srcpath{sparse\_kron\_reduction.hpp} —— 各类降阶；
  \item \srcpath{include/hacdcpf/graph/reduction\_plan.hpp}、
        \srcpath{reduction\_mapping.hpp}、\srcpath{result\_recovery.hpp} —— 降阶计划、映射与状态恢复；
  \item \srcpath{docs/developer/graph\_runtime\_contract.md} —— 既有运行时契约。
\end{itemize}

\subsection{单位制与索引空间}
支路参数 $r/x/b/\text{tap}$ 为标幺值（pu）；本模块严格区分三类 ID：组件
\fld{comp\_index}（对外稳定）、\fld{edge\_id}（0-based 位置索引）与图 node/edge 临时索引。
AC 与 DC 母线各自经 \fld{ac\_bus\_id\_to\_node\_idx}/\fld{dc\_bus\_id\_to\_node\_idx} 映射到
node 下标，恢复结果 \texttt{FullNetworkVoltages} 按 \fld{ac\_bus\_voltage}/\fld{dc\_bus\_voltage}
分列。零阻抗判定阈值为 \fld{zero\_impedance\_threshold}（默认 $10^{-8}$）。

\subsection{符号约定}
$N$ 节点数、$E$ 边数、$p$ 连通分量数、$\mu$ 圈复杂度；$Y$ 母线导纳矩阵，按保留集 $\alpha$
与消去集 $\beta$ 分块；$Z_{ij}$ 支路阻抗；$V_i$ 母线电压相量。降阶映射统一记为
\texttt{ReductionMapping}，状态恢复算子记 \srcpath{recover\_*}。

\subsection{诚实口径}
每类降阶都记录一份映射以支持 \srcpath{recover\_*} 状态恢复：Kron 仅对无源/线性内部节点严格
精确，恒功率内部节点为工作点相关的近似；填充比守卫 \fld{max\_fill\_ratio} 与奇异
$Y_{\beta\beta}$ 主元守卫可拒绝降阶；悬垂折叠为近似（取 $|V_{\text{parent}}|=1.0$ pu，迭代恢复
不单独报告收敛标志）。旧式扁平映射（\fld{bus\_id\_to\_node\_idx}、\texttt{IslandInfo::bus\_ids}
等）在 AC/DC 同号 ID 时含糊，须改用 domain-qualified 映射。

\section{模块定位与工程应用场景}
\subsection{功能定位}
图模块是投影层与装配层之间的\textbf{拓扑基础设施}：它把工程语义系统建为无向图，提供
连通性/辐射性/桥/割点等拓扑判定，并在装配 $Y_{bus}$ 之前对网络做保序降阶（开关收缩、
串联/悬垂合并、Kron/Schur 补），同时为每类降阶保留可逆映射以在求解后恢复被消去母线的
电压与支路潮流。它是潮流、OPF、重构、可靠性、弹性等上层模块共享的公共层。

\subsection{工程应用场景}
典型场景：大规模系统求解前的规模压缩（Kron 消去无源内部节点，降低 $Y_{bus}$ 维数）；
含大量闭合开关/断路器的配电网建模（零阻抗收缩为超节点）；孤岛与辐射性检测（重构、恢复、
黑启动前置判定）；桥/割点识别（可靠性 N-1 与结构脆弱性分析）。降阶后的解经映射恢复回
原始编号空间，对上层完全透明。

\section{输入、输出与边界条件}
\subsection{输入}
\texttt{HybridPowerSystem}（含 AC/DC 母线、支路、开关、断路器、VSC、DC-DC 耦合）；
零阻抗阈值 \fld{zero\_impedance\_threshold}；降阶选项 \texttt{GraphReductionOptions}
（各降阶开关、\fld{max\_fill\_ratio}、\fld{mode}、\fld{preserve\_*} 强制保留集）。

\subsection{输出}
\texttt{PowerSystemGraph}（节点/边/邻接表 + domain-qualified 映射）；\texttt{TopologyReport}
（孤岛列表、连通/辐射标志、回路数、桥边、割点、基本圈、诊断）；降阶后的等值网络 +
\texttt{ReductionMapping}；状态恢复输出 \texttt{FullNetworkVoltages}（AC/DC 分列）。

\subsection{边界条件}
零阻抗合并要求两端电压基一致（容差内）；多平衡节点合并默认拒绝（除非
\fld{allow\_multi\_slack\_merge}）；Kron 要求 $Y_{\beta\beta}$ 非奇异且填充比不超限；
串联降阶要求中间节点无源且度为 2；被 \fld{preserve\_*} 标记的母线/支路不参与任何降阶。

\section{变量、参数与符号定义}
\subsection{关键数据结构}
\compmeta{PowerSystemGraph}{include/hacdcpf/graph/power\_system\_graph.hpp}
主要字段：\fld{nodes}（\texttt{GraphNode}）、\fld{edges}（\texttt{GraphEdge}：\fld{comp\_index}、
\fld{edge\_id}、\fld{r\_pu}/\fld{x\_pu}/\fld{b\_pu}/\fld{tap}、\fld{is\_zero\_impedance}）、
\fld{ac\_bus\_id\_to\_node\_idx}、\fld{dc\_bus\_id\_to\_node\_idx}、\fld{adj}。

\compmeta{TopologyReport}{include/hacdcpf/graph/topology\_analysis.hpp}
主要字段：\fld{islands}、\fld{is\_connected}、\fld{is\_radial}、\fld{cycle\_count}、
\fld{bridge\_edge\_ids}、\fld{cut\_vertices}、\fld{fundamental\_cycles}、\fld{diagnostics}。

\compmeta{GraphReductionOptions}{include/hacdcpf/graph/reduction\_plan.hpp}
主要字段：\fld{enable\_switch\_contraction}/\fld{enable\_series\_reduction}/
\fld{enable\_pendant\_reduction}/\fld{enable\_kron\_reduction}、\fld{zero\_impedance\_threshold}、
\fld{max\_fill\_ratio}、\fld{mode}（\texttt{Safe}/\texttt{Moderate}/\texttt{Aggressive}）。
降阶结果映射为 \texttt{ReductionMapping}，恢复输出为 \texttt{FullNetworkVoltages}。

\subsection{参数字典}
\begin{paramtable}
\fld{zero\_impedance\_threshold} & — & pu & 小正数 & $10^{-8}$ & 零阻抗/闭合开关判定阈值\\
\fld{max\_fill\_ratio} & — & — & $\ge 1$ & $2.0$ & Kron 填充比上限（超限拒绝）\\
\fld{enable\_switch\_contraction} & — & 布尔 & — & \texttt{true} & 是否启用开关/零阻抗收缩\\
\fld{enable\_series\_reduction} & — & 布尔 & — & — & 是否启用串联降阶\\
\fld{enable\_pendant\_reduction} & — & 布尔 & — & — & 是否启用悬垂折叠（近似）\\
\fld{enable\_kron\_reduction} & — & 布尔 & — & — & 是否启用 Kron/Schur 补降阶\\
\fld{mode} & — & — & Safe/Moderate/Aggressive & \texttt{Safe} & 降阶激进程度\\
\fld{allow\_multi\_slack\_merge} & — & 布尔 & — & \texttt{false} & 是否允许多平衡节点合并\\
\end{paramtable}
\srcpath{build\_power\_system\_graph} 的 \fld{zero\_impedance\_threshold} 参数默认 $10^{-8}$；
\fld{preserve\_*} 系列标志用于强制保留指定母线/支路不被降阶。

% ── 通用理论层（修订版 2 新增，专著级深度）──
% theory_graph_foundations.tex —— 图模块 通用理论：图论基础、网络矩阵与降阶理论
% 本文件为理论层章节，遵循 docs/modules/THEORY_WRITING_GUIDE.md。

\section{图与网络降阶的数学基础}\label{sec:gr-th-foundations}

\begin{theorynote}
本章为通用数学理论，按图论与稀疏矩阵计算的标准文献体系撰写，覆盖：图论基本概念
（图、子图、度、路径、连通分量、邻接表示）；电网络的矩阵表示（节点-支路关联矩阵、
KCL/KVL 的矩阵形式、$Y_{bus}=AY_bA^{\mathsf T}$ 的推导及其代数性质）；树与余树理论
（基本回路/割集、辐射性的定理化刻画）；DFS/BFS 与 Tarjan 桥/割点算法的正确性论证；
Kron 降阶一般理论（Schur 补与高斯消元的等价、精确性条件、含注入的 Ward 等值）；
稀疏性理论（fill-in 与消去顺序、最小度/AMD/嵌套分割）；以及多级降阶的映射复合代数、
星-三角/串并联等值与近似等值的误差讨论。本章公式不要求逐式对应代码；与平台实现
（\srcpath{src/graph/}、\srcpath{include/hacdcpf/graph/}）的对应关系见本章末
"与实现的对应"表，超出实现的内容以"未实现扩展说明"框就地标记。
\end{theorynote}

\subsection{记号与约定}
\begin{itemize}
  \item $G=(V,E)$：无向图，$N=|V|$ 节点数，$E=|E|$ 边数，$p$ 连通分量数；
        允许并联边（电力网络中双回线常见），不允许自环；
  \item $\mu=E-N+p$：圈复杂度（cyclomatic number，独立回路数）；
  \item $A\in\{0,\pm1\}^{N\times E}$：节点-支路关联矩阵；$B$：基本回路矩阵；
        $Y_b=\operatorname{diag}(y_e)$：支路原始导纳矩阵；
  \item $Y\in\mathbb C^{N\times N}$：节点导纳矩阵；$\alpha$ 保留集、$\beta$ 消去集，
        分块记 $Y_{\alpha\alpha},Y_{\alpha\beta},Y_{\beta\alpha},Y_{\beta\beta}$；
  \item $V_i$：节点电压相量；$I_i$：节点注入电流；$i_e$：支路电流（沿边定向）；
  \item $disc[\cdot]$、$low[\cdot]$：DFS 发现时间与 low-link 值；
  \item $\operatorname{nnz}(\cdot)$：矩阵非零元个数；$r_{fill}$：填充比。
\end{itemize}
与本手册其余部分冲突时以本章声明为准。

\subsection{图论基础形式化}\label{sec:gr-th-basics}

\begin{definition}[图、子图与度]
\emph{无向图} $G=(V,E)$ 由节点集 $V$ 与边集 $E$ 组成，每条边 $e\in E$ 关联无序节点对
$\{u,v\}$。\emph{子图} $G'=(V',E')$ 满足 $V'\subseteq V$、$E'\subseteq E$；若
$V'=V$ 则称\emph{生成子图}。节点 $v$ 的\emph{度} $\deg(v)$ 为与其关联的边数。
电力网络建模中，母线对应节点，支路/开关/换流器耦合对应边；并联支路使图成为
多重图，本章结论对多重图同样成立。
\end{definition}

\begin{lemma}[握手引理]\label{sec:gr-th-handshake}
$\sum_{v\in V}\deg(v)=2E$。
\end{lemma}
\begin{proof}
对关联关系集合 $\mathcal I=\{(v,e):v\in e\}$ 双向计数：按节点求和得
$|\mathcal I|=\sum_v\deg(v)$；按边求和，每条边恰关联两个端点，得
$|\mathcal I|=2E$。两者相等。
\end{proof}

\begin{definition}[路径、连通与连通分量]
\emph{路径}是节点序列 $v_0,v_1,\dots,v_k$，相邻节点间有边相连；$v_0=v_k$ 且边互不
重复时称\emph{回路}（圈）。若任意两节点间存在路径，称图\emph{连通}。连通关系是
$V$ 上的等价关系，其等价类导出的极大连通子图称为\emph{连通分量}（电力语境下的
\emph{孤岛}），分量数记为 $p$。
\end{definition}

\begin{definition}[邻接矩阵与邻接表]
\emph{邻接矩阵} $\mathcal A(G)\in\{0,1\}^{N\times N}$ 定义为
$\mathcal A_{ij}=$（连接 $i,j$ 的边数）。\emph{邻接表}为每个节点存储其
（边号, 邻居）对列表。
\end{definition}

\begin{remark}[表示法的选择]
邻接矩阵占用 $O(N^2)$ 空间，适合稠密小图的代数运算；邻接表占用 $O(N+E)$ 空间，
支持 $O(\deg(v))$ 的邻居枚举，是 DFS/BFS 等遍历算法达到线性复杂度的前提。
电力网络极度稀疏（每节点平均度通常为 $2$–$5$），邻接表是唯一合理的遍历表示；
而邻接矩阵的代数推广——节点导纳矩阵——则以稀疏格式参与数值求解。
\end{remark}

\subsection{电网络的矩阵表示}\label{sec:gr-th-network}

\begin{definition}[节点-支路关联矩阵]
为每条边 $e=\{i,j\}$ 任意指定定向 $i\to j$，定义
$A\in\{0,\pm1\}^{N\times E}$：
\begin{equation}
A_{ve}=\begin{cases}
+1,& e \text{ 以 } v \text{ 为起点},\\
-1,& e \text{ 以 } v \text{ 为终点},\\
0,& \text{其余}.
\end{cases}
\end{equation}
定向是计算约定，不影响任何物理结论（支路电流符号随之翻转）。
\end{definition}

\begin{proposition}[关联矩阵的秩]\label{sec:gr-th-rankA}
设图有 $p$ 个连通分量，则 $\operatorname{rank}A=N-p$（实数域上）。
\end{proposition}
\begin{proof}
每列恰含一个 $+1$ 与一个 $-1$，故同一分量内全部行向量之和为零，得
$\operatorname{rank}A\le N-p$。下设图连通（$p=1$），证 $\operatorname{rank}A\ge N-1$：
取一棵生成树（第~\ref{sec:gr-th-tree}~节证明其存在），对其 $N-1$ 条树边对应的列
做归纳。树叶（度为 1 的树节点）对应的行在其唯一关联树边列上取 $\pm1$、其余树边列
上取 $0$；删去该叶与其树边后剩余仍是树。据此可从树叶向根逐次做行消元，将树边列
组成的 $N\times(N-1)$ 子矩阵化为阶梯形，含 $N-1$ 个主元，故秩至少 $N-1$。
各分量独立成立，合并即得 $\operatorname{rank}A=N-p$。
\end{proof}

\begin{definition}[KCL/KVL 的矩阵形式]
记支路电流向量 $i_b\in\mathbb C^{E}$（沿边定向）、节点注入电流向量
$I\in\mathbb C^{N}$（注入网络为正）、节点对地电压向量 $V\in\mathbb C^{N}$，
支路电压向量 $v_b\in\mathbb C^{E}$（沿定向的电压降）。则
\begin{equation}
\text{KCL：}\quad I=A\,i_b;\qquad
\text{KVL：}\quad v_b=A^{\mathsf T}V.
\end{equation}
第一式逐行读：节点 $v$ 的注入等于以其为起点的支路电流之和减去以其为终点的支路
电流之和；第二式逐列读：支路电压降等于两端节点电压之差，对任意回路求和为零。
\end{definition}

\begin{definition}[支路原始导纳]
忽略支路间互感耦合，每条支路以集中参数 $\pi$ 等值（阻抗 $z_e=r_e+\mathrm jx_e$、
对地充电纳入节点并联项），支路导纳 $y_e=1/z_e$，原始导纳矩阵
$Y_b=\operatorname{diag}(y_1,\dots,y_E)$。
\end{definition}

\begin{theorem}[节点导纳矩阵的关联矩阵分解]\label{sec:gr-th-ybus}
在不计互感、支路电流满足 $i_b=Y_bv_b$ 的网络中，
\begin{equation}
Y_{bus}=A\,Y_b\,A^{\mathsf T},
\end{equation}
且其元素为
\begin{equation}
[Y_{bus}]_{ij}=-\sum_{e\in\,i\!\sim\!j}y_e\ \ (i\ne j),
\qquad
[Y_{bus}]_{ii}=\sum_{e\ni i}y_e+y_i^{\mathrm{sh}},
\end{equation}
其中 $e\in i\!\sim\!j$ 遍历连接 $i,j$ 的全部（并联）边，$y_i^{\mathrm{sh}}$ 为节点
$i$ 的对地并联导纳（充电、负荷导纳、接地支路之和，作为对 $A Y_b A^{\mathsf T}$
对角元的追加）。
\end{theorem}
\begin{proof}
由 KCL/KVL 与支路关系：$I=Ai_b=AY_bv_b=AY_bA^{\mathsf T}V$，故
$Y_{bus}=AY_bA^{\mathsf T}$。元素计算：
$[AY_bA^{\mathsf T}]_{ij}=\sum_e y_eA_{ie}A_{je}$。$i\ne j$ 时乘积 $A_{ie}A_{je}$
非零当且仅当 $e$ 同时关联 $i$ 与 $j$，此时 $e$ 以一端为起点、另一端为终点，
乘积为 $-1$，故 $[Y_{bus}]_{ij}=-\sum_{e\in i\sim j}y_e$；$i=j$ 时
$A_{ie}^2$ 对每条关联边取 $1$，故对角元为关联支路导纳之和，再计入节点对地
并联项 $y_i^{\mathrm{sh}}$。
\end{proof}

\begin{remark}
上式给出 $Y_{bus}$ 的稀疏结构：$[Y_{bus}]_{ij}\ne0$（$i\ne j$）当且仅当
$i,j$ 相邻——$Y_{bus}$ 的零模式恰为图邻接结构的代数镜像，这是第~\ref{sec:gr-th-sparsity}~节
稀疏性理论的出发点。含移相变压器或不对称耦合时 $Y_{bus}$ 不再对称，分解式须推广；
本章默认对称情形。
\end{remark}

\begin{theorem}[接地导纳矩阵的正定性]\label{sec:gr-th-ybus-pd}
设网络连通、每条支路电导 $g_e=\Re y_e>0$、存在对地通路（至少一个节点的
$y_i^{\mathrm{sh}}$ 有正实部，或指定参考节点接地）。删去参考节点行列后所得的
\emph{接地导纳矩阵} $\bar Y$ 满足：对任意 $V\ne0$，
$\Re\bigl(V^{\mathsf H}\bar YV\bigr)>0$，特别地 $\bar Y$ 非奇异。
\end{theorem}
\begin{proof}
由分解式，
\begin{equation}
\Re\bigl(V^{\mathsf H}Y_{bus}V\bigr)
=\sum_{e=\{i,j\}}g_e\,|V_i-V_j|^2+\sum_i \Re\bigl(y_i^{\mathrm{sh}}\bigr)|V_i|^2\ge0.
\end{equation}
若该式为零，则每条边上 $|V_i-V_j|=0$，连通性迫使所有节点电压相等为某常数 $c$；
对地通路（$\Re y_i^{\mathrm{sh}}>0$ 或接地约束 $V_{\mathrm{ref}}=0$）迫使 $c=0$。
故半正定型仅在 $V=0$ 时取零，$\bar Y$ 实部正定、非奇异。
\end{proof}

\begin{definition}[基本回路矩阵]
给定生成树（或生成森林），每条\emph{连支}（非树边）与树上的唯一路径构成一个
\emph{基本回路}（第~\ref{sec:gr-th-tree}~节）。基本回路矩阵
$B\in\{0,\pm1\}^{\mu\times E}$ 的每行对应一个基本回路，沿回路取向在所含边上取
$\pm1$。KVL 的回路形式为
\begin{equation}
B\,v_b=0.
\end{equation}
\end{definition}

\subsection{树、余树与辐射性}\label{sec:gr-th-tree}

\begin{definition}[树、生成树与余树]
连通无圈图称为\emph{树}。图 $G$ 的\emph{生成树}是包含全部节点的树形生成子图；
生成树之外的边称\emph{连支}（弦），其集合称\emph{余树}。对 $p$ 个分量的图，
逐分量取生成树得\emph{生成森林}。
\end{definition}

\begin{theorem}[树的等价刻画]\label{sec:gr-th-tree-equiv}
对含 $N$ 个节点的图 $G$，以下命题等价：
\begin{enumerate}
  \item $G$ 是树（连通且无圈）；
  \item $G$ 连通且 $E=N-1$；
  \item $G$ 无圈且 $E=N-1$；
  \item 任意两节点间存在\emph{唯一}路径。
\end{enumerate}
\end{theorem}
\begin{proof}
(i)$\Rightarrow$(ii)：对 $N$ 归纳。$N=1$ 时 $E=0$ 成立。$N\ge2$ 时树必有叶
（否则沿边游走必重复节点成圈），删去一叶及其关联边得 $N-1$ 节点的树，由归纳假设
其边数为 $N-2$，故 $E=N-1$；连通性显然。

(ii)$\Rightarrow$(i)：设 $G$ 连通且 $E=N-1$。若 $G$ 含圈，删去圈上一边不破坏
连通性（圈上两点仍有绕行路径），重复至无圈，得连通无圈子图 $G'$ 即树，
由已证方向其边数为 $N-1$；但每删一边边数严格减少，与 $E=N-1$ 矛盾，故原图本无圈。

(i)$\Rightarrow$(iv)：连通保证路径存在；若两点间有两条不同路径，其对称差含圈，
矛盾。

(iv)$\Rightarrow$(i)：路径存在即连通；若含圈，圈上两点间有两条路径，矛盾。

(iii) 与 (i)(ii) 的互推由同一计数论证给出：无圈图逐分量归纳得 $E=N-p$，
$E=N-1$ 当且仅当 $p=1$（连通）。
\end{proof}

\begin{corollary}[辐射性的代数判据]\label{sec:gr-th-radial}
配电网络\emph{辐射状}（radial）当且仅当其图连通且 $E=N-1$；等价地，连通且
圈复杂度 $\mu=0$。这为重构、恢复与黑启动的拓扑约束提供了 $O(N+E)$ 可判的
充要条件。
\end{corollary}

\begin{theorem}[圈空间维数与基本回路基]\label{sec:gr-th-cyclebasis}
$\mu=E-N+p$ 恰为图的\emph{圈空间}维数（边子集在 $\mathrm{GF}(2)$ 上按对称差
构成的向量空间中，由全体回路生成的子空间之维数）。任取生成森林，其 $\mu$ 条
连支各确定一个基本回路；这组基本回路构成圈空间的一组基。
\end{theorem}
\begin{proof}
生成森林含 $N-p$ 条边（每分量 $N_c-1$ 条，对分量求和），故连支数为
$E-(N-p)=\mu$。每个基本回路恰含一条其他基本回路不含的连支，故它们线性无关
（任何非空对称差必含某连支而非空集）。再证生成性：任一回路 $C$，取其所含连支
对应的基本回路之对称差 $S$；$C\triangle S$ 是无圈子集且只含树边，而无圈子集
中若含边则其导出子图有圈——唯一可能是空集，故 $C=S$。于是 $\mu$ 个基本回路
构成基，维数为 $\mu$。
\end{proof}

\begin{definition}[基本割集]
删去生成树的任一树边将树分为两部分，跨这两部分的全部原图边构成一个
\emph{基本割集}。$N-p$ 条树边对应 $N-p$ 个基本割集，它们构成割空间
（边割按对称差生成的空间）的一组基，维数 $N-p$。圈空间与割空间互为正交补，
维数之和 $\mu+(N-p)=E$。
\end{definition}

\begin{remark}
基本回路/割集理论是经典回路法与割集法潮流方程的基础，也直接支撑拓扑报告中的
回路计数 $\fld{cycle\_count}=\max(0,E-N+p)$ 与基本圈枚举。
\end{remark}

\subsection{DFS、BFS 与桥/割点算法}\label{sec:gr-th-dfs}

\begin{definition}[DFS 树、发现时间与回边]
深度优先搜索（DFS）从未访问节点出发递归深入，首次到达节点 $v$ 的时刻记
$disc[v]$（时间戳从 $1$ 递增）。DFS 的递归结构生成\emph{DFS 生成森林}：
首次发现 $v$ 时所经边为\emph{树边}，其余为\emph{非树边}。
\end{definition}

\begin{lemma}[无向图无横叉边]\label{sec:gr-th-nocross}
无向图的 DFS 中，每条非树边必连接一对祖先-后裔节点（\emph{回边}），不存在
连接两棵不相干子树的横叉边。
\end{lemma}
\begin{proof}
设边 $\{u,v\}$ 存在且 $disc[u]<disc[v]$。DFS 在离开 $u$ 之前必考察
$\{u,v\}$；若此时 $v$ 未访问，则 $\{u,v\}$ 成为树边；若 $v$ 已访问，则
$disc[u]<disc[v]$ 且 $v$ 的递归必已结束，由栈结构 $v$ 是 $u$ 的真后裔
（否则 $v$ 先开始且先结束，则考察该边时 $u$ 尚未入栈，矛盾）。故非树边
连接祖先与后裔。
\end{proof}

\begin{definition}[low-link 值]
\begin{equation}
low[v]=\min\Bigl(disc[v],\ \min_{\substack{u\in\mathrm{subtree}(v)\\
\{u,w\}\,\text{回边}}}disc[w]\Bigr),
\end{equation}
即 $v$ 的子树经树边向下、再经至多一条回边向上所能触及的最小发现时间。
DFS 中可递推计算：
$low[v]=\min\bigl(disc[v],\ \min_{w:\,\text{回边}}disc[w],\ \min_{u\,\text{子节点}}low[u]\bigr)$。
\end{definition}

\begin{lemma}[桥判据]\label{sec:gr-th-bridge}
设 $(u,v)$ 为树边，$u$ 为 $v$ 的父节点。则 $\{u,v\}$ 是桥（删去后图不再连通）
当且仅当 $low[v]>disc[u]$。
\end{lemma}
\begin{proof}
（$\Leftarrow$）若 $low[v]>disc[u]$，则 $v$ 的子树内不存在指向 $u$ 或其祖先的
回边；由引理~\ref{sec:gr-th-nocross}，$v$ 子树与图其余部分之间除 $\{u,v\}$ 外
无其他边，删去 $\{u,v\}$ 后 $v$ 子树脱离，故为桥。

（$\Rightarrow$）若 $\{u,v\}$ 是桥，则 $v$ 子树与其余部分之间无其他边，特别地
无回边指向发现时间 $\le disc[u]$ 的节点，故 $low[v]\ge disc[v]>disc[u]$。
\end{proof}

\begin{lemma}[割点判据]\label{sec:gr-th-ap}
设 $u$ 非 DFS 根。$u$ 是割点（删去后图连通分量数增加）当且仅当存在子节点 $v$
使 $low[v]\ge disc[u]$。DFS 根是割点当且仅当其在 DFS 树中有至少两个子节点。
\end{lemma}
\begin{proof}
非根情形。（$\Leftarrow$）若子节点 $v$ 满足 $low[v]\ge disc[u]$，则 $v$ 子树
无回边绕过 $u$ 到达其真祖先；删去 $u$ 后 $v$ 子树与 $u$ 的祖先部分断开，分量数
增加。（$\Rightarrow$）若 $u$ 是割点，删去 $u$ 后图分裂，其中必有一块由 $u$ 的
某棵子树 $v$ 整体构成（DFS 树结构保证跨子树无边，引理~\ref{sec:gr-th-nocross}），
该子树无回边指向 $disc$ 小于 $disc[u]$ 的节点，即 $low[v]\ge disc[u]$。

根情形：根无祖先，其删除后各 DFS 子树两两之间无边（同引理），故根为割点当且仅当
子树数 $\ge2$。
\end{proof}

\begin{theorem}[Tarjan 桥/割点算法]\label{sec:gr-th-tarjan}
单趟 DFS 依据引理~\ref{sec:gr-th-bridge}、\ref{sec:gr-th-ap} 的判据，可在
$O(N+E)$ 时间与 $O(N)$ 附加空间内求出全部桥与割点~\cite{grf:tarjan}。
\end{theorem}
\begin{proof}
正确性：$disc$、$low$ 在 DFS 回溯时按递推式 $O(\deg(v))$ 逐节点计算，low-link
递推的正确性由定义直接归纳（子节点的 $low$ 已先算好）；判定条件即两条引理，
穷举所有树边与非根节点，故输出恰为全部桥与割点。

复杂度：DFS 每节点入栈一次、每条边在两端正向各考察一次，$low$ 更新与判据测试
均为 $O(1)$，总计 $O(N+E)$；存储 $disc/low/$栈为 $O(N)$。
\end{proof}

\begin{remark}[BFS 的角色]
广度优先搜索（BFS）按层展开，给出到源点的最少边数距离与一棵矮生成树，适合
连通分量枚举与基本回路的路径回溯（沿父指针上升至最近公共祖先）。DFS 则天然给出
祖先-后裔结构，是 low-link 类算法的载体。两者同为 $O(N+E)$，在平台中分别承担
孤岛枚举/基本圈构造与桥/割点检测。
\end{remark}

\subsection{Kron 降阶与 Ward 等值}\label{sec:gr-th-kron}

\begin{definition}[分块与 Schur 补]
将节点集划分为保留集 $\alpha$ 与消去集 $\beta$，网络方程 $YV=I$ 分块为
\begin{equation}
\begin{bmatrix}Y_{\alpha\alpha}&Y_{\alpha\beta}\\[2pt]
Y_{\beta\alpha}&Y_{\beta\beta}\end{bmatrix}
\begin{bmatrix}V_\alpha\\V_\beta\end{bmatrix}
=
\begin{bmatrix}I_\alpha\\I_\beta\end{bmatrix}.
\end{equation}
设 $Y_{\beta\beta}$ 非奇异，定义 \emph{Schur 补}
\begin{equation}
Y_{red}=Y_{\alpha\alpha}-Y_{\alpha\beta}Y_{\beta\beta}^{-1}Y_{\beta\alpha},
\qquad\text{记 }K=Y_{\beta\beta}^{-1}Y_{\beta\alpha}.
\end{equation}
\end{definition}

\begin{theorem}[Schur 补与高斯消元的等价]\label{sec:gr-th-schur-gauss}
对分块方程组做一步块高斯消元（以 $Y_{\beta\beta}$ 为主元块消去 $V_\beta$）
所得关于 $V_\alpha$ 的方程恰为
\begin{equation}
Y_{red}\,V_\alpha=I_\alpha-Y_{\alpha\beta}Y_{\beta\beta}^{-1}I_\beta,
\qquad
V_\beta=Y_{\beta\beta}^{-1}\bigl(I_\beta-Y_{\beta\alpha}V_\alpha\bigr).
\end{equation}
反之，由上式解出的 $(V_\alpha,V_\beta)$ 必满足原方程组。即消去是\emph{精确的
同解变换}，不产生任何近似。
\end{theorem}
\begin{proof}
第二块行左乘 $Y_{\alpha\beta}Y_{\beta\beta}^{-1}$ 后从第一块行减去：
\begin{equation}
\bigl(Y_{\alpha\alpha}-Y_{\alpha\beta}Y_{\beta\beta}^{-1}Y_{\beta\alpha}\bigr)V_\alpha
=I_\alpha-Y_{\alpha\beta}Y_{\beta\beta}^{-1}I_\beta,
\end{equation}
即所断言的约化方程；第二块行移项即得 $V_\beta$ 的恢复式。反之，由约化方程解出
$V_\alpha$、由恢复式定义 $V_\beta$，回代验证两块行均成立，故同解。
\end{proof}

\begin{lemma}[消去块的非奇异充分条件]\label{sec:gr-th-ybb}
若 $Y$ 严格对角占优（按行，复数情形按模），则其任意主子阵——特别地
$Y_{\beta\beta}$——非奇异。
\end{lemma}
\begin{proof}
主子阵的行仅删去若干非对角元，对角元不变，故严格对角占优性保持。由 Gershgorin
圆盘定理，严格对角占优矩阵的所有特征值模不小于
$\min_i\bigl(|Y_{ii}|-\sum_{j\ne i}|Y_{ij}|\bigr)>0$，故非奇异。
电力网络中节点并联项（充电、负荷导纳）常使 $Y$ 满足（不可约）对角占优，
但对角占优是充分而非必要条件，工程上仍需显式主元守卫。
\end{proof}

\begin{theorem}[Kron 降阶的精确性条件]\label{sec:gr-th-kron-exact}
约化导纳模型 $Y_{red}V_\alpha=I_\alpha$ 对\emph{任意}保留集注入 $I_\alpha$
都与原网络同解，当且仅当消去节点无注入（$I_\beta=0$，无源节点）。此时
\begin{equation}
Y_{red}=Y_{\alpha\alpha}-Y_{\alpha\beta}Y_{\beta\beta}^{-1}Y_{\beta\alpha},
\qquad
V_\beta=-K\,V_\alpha=-Y_{\beta\beta}^{-1}Y_{\beta\alpha}V_\alpha,
\end{equation}
降阶对保留节点电压、注入与端口潮流\emph{逐量精确}（Kron 降阶~\cite{grf:kron}）。
\end{theorem}
\begin{proof}
由定理~\ref{sec:gr-th-schur-gauss}，含注入修正项的约化方程对任意 $I_\alpha$
同解；令 $I_\beta=0$ 即退化为 $Y_{red}V_\alpha=I_\alpha$，恢复式退化为
$V_\beta=-KV_\alpha$。必要性：若 $I_\beta\ne0$ 而仍用无修正的约化模型，则两模型
解相差 $Y_{red}^{-1}Y_{\alpha\beta}Y_{\beta\beta}^{-1}I_\beta$；该偏差向量
对一般 $I_\beta$ 非零（例如取 $I_\beta$ 沿 $Y_{\beta\beta}$ 的某列方向），
故精确性要求 $I_\beta=0$。
\end{proof}

\begin{theorem}[含注入的 Ward 等值]\label{sec:gr-th-ward}
设消去节点带\emph{恒定电流}注入 $I_\beta$（与 $V$ 无关），则精确的等值网络
由约化导纳矩阵与修正注入组成（Ward 等值~\cite{grf:ward}）：
\begin{equation}
Y_{red}=Y_{\alpha\alpha}-Y_{\alpha\beta}Y_{\beta\beta}^{-1}Y_{\beta\alpha},
\qquad
I_{red}=I_\alpha-Y_{\alpha\beta}Y_{\beta\beta}^{-1}I_\beta,
\end{equation}
求解 $Y_{red}V_\alpha=I_{red}$ 后，消去节点电压按下式精确恢复：
\begin{equation}
V_\beta=Y_{\beta\beta}^{-1}\bigl(I_\beta-Y_{\beta\alpha}V_\alpha\bigr)
=-K\,V_\alpha+Y_{\beta\beta}^{-1}I_\beta.
\end{equation}
\end{theorem}
\begin{proof}
即定理~\ref{sec:gr-th-schur-gauss} 的同解变换按 $I_{red}$ 定义重写：
$I_{red}$ 恰为约化方程右端；恢复式即第二块行移项。全程仅为线性代数恒等变换，
对恒定 $I_\beta$ 无近似。
\end{proof}

\begin{remark}[恒功率注入导致的近似]
交流潮流中负荷为恒功率注入 $I_j=(S_j/V_j)^{*}$，依赖被消去的电压本身，
不再满足 Ward 等值的线性前提。此时 Kron/Ward 消去相当于在基准工作点处把
恒功率负荷冻结为恒电流或恒导纳的近似，其误差随扰动后运行点偏离基准点而增长；
近似等值的保真度本质上是\emph{工作点相关}的，必须结合灵敏度或误差界评估。
\end{remark}

\begin{gapnote}
当前稠密 Kron 入口 \srcpath{apply\_kron\_reduction\_with\_injection} 已实现
$I_{red}$ 的计算，但恢复 API \srcpath{recover\_eliminated\_voltages} 只返回
$V_\beta=-KV_\alpha$，未含定理~\ref{sec:gr-th-ward} 中的
$+Y_{\beta\beta}^{-1}I_\beta$ 修正项；带非零 $I_\beta$ 的完整电压恢复在现有
接口下不可达。恒功率内部节点 Kron 近似的工作点相关误差界估计亦未实现
（见手册"后续改进方向"）。
\end{gapnote}

\begin{gapnote}
更精细的外部系统等值——REI 等值（按设备类别聚合注入的辐射等值独立网）、
Dimo 等值等——未在图模块实现；当前仅覆盖 Ward/Kron 一族。
\end{gapnote}

\subsection{稀疏性与消去顺序}\label{sec:gr-th-sparsity}

\begin{definition}[fill-in 与消去图]
对稀疏矩阵做消元时，原本为零的位置因消去运算变为非零，称为\emph{填入}
（fill-in）。把矩阵的零模式看作图（非零元 $Y_{ij}$ 对应边 $\{i,j\}$），
消去过程对应\emph{消去图}序列 $G^{(0)}=G,G^{(1)},\dots$：消去节点 $k$ 后
从图中删除 $k$ 及其关联边。
\end{definition}

\begin{theorem}[消去的成团规则；fill-path 刻画]\label{sec:gr-th-fill}
消去节点 $k$ 时，其当前邻居在消去图中两两相连成为团（clique）：
\begin{equation}
Y_{ij}\leftarrow Y_{ij}-\frac{Y_{ik}Y_{kj}}{Y_{kk}},\qquad i,j\in N(k),
\end{equation}
只要 $Y_{ik}Y_{kj}\ne0$ 且原本 $Y_{ij}=0$，位置 $(i,j)$ 即被填入。
更一般地（Rose~\cite{grf:rose}）：节点对 $(i,j)$ 在最终三角因子中为非零，
当且仅当原图中存在一条从 $i$ 到 $j$ 的路径，其全部内部节点在消去顺序中
都先于 $i$ 与 $j$ 被消去。
\end{theorem}
\begin{proof}
成团规则：消去 $k$ 即对剩余方程做 Schur 补更新，$(i,j)$ 位置增量为
$-Y_{ik}Y_{kj}/Y_{kk}$，$i,j$ 同为 $k$ 的邻居时增量一般非零，故邻居两两相连。
fill-path 刻画对消去步数归纳：基础步即成团规则（路径 $i$--$k$--$j$，内部节点
仅 $k$ 且先被消去）。归纳步：若因子中 $(i,j)$ 非零，则或原图有边（空内部路径），
或某步消去 $k$ 时填入——由归纳假设 $i$ 到 $k$、$k$ 到 $j$ 各有内部节点先于
$\min(i,j)$ 被消去的路径，拼接（经 $k$，且 $k$ 先于 $i,j$ 被消去）即得所求路径。
反向由同一归纳成立。
\end{proof}

\begin{remark}[消去顺序的决定性影响]
同一矩阵不同消去顺序产生的填入量可相差数量级。电力网络 $Y_{bus}$ 极端稀疏，
消去顺序实质上决定了 LU 分解的内存与时间。最小填入消去顺序问题是 NP 完全的
（Yannakakis~\cite{grf:yannakakis}），工程上依赖启发式。
\end{remark}

\begin{definition}[最小度、AMD 与嵌套分割]
\emph{最小度}（Minimum Degree，Tinney 第二方案~\cite{grf:tinney}）：每步在消去图中
选取当前度最小的节点消去——其邻居成团至多引入 $\binom{d}{2}$ 条新边，贪心压低
单步填入。\emph{近似最小度}（AMD）以可廉价更新的度上界代替精确度，摊还复杂度
更低，是大规模稀疏库（SuiteSparse 等）的默认排序。\emph{嵌套分割}（Nested
Dissection~\cite{grf:georgeliu}）：递归地用顶点分隔集把图劈为两半，最后消去分隔集；
对 $n\times n$ 网格类图可证填入 $O(N\log N)$、分解工作量 $O(N^{3/2})$
（$N=n^2$），远优于自然顺序的 $O(N^2)$ 填入。
\end{definition}

\begin{definition}[填充比]
降阶（或消元）前后非零元之比
\begin{equation}
r_{fill}=\frac{\operatorname{nnz}(Y_{red})}{\operatorname{nnz}(Y_{full})}
\end{equation}
度量消去集选择对稀疏性的破坏。消去高度节点会制造大团，使 $r_{fill}\gg1$，
"降阶"反而增加求解成本；因此工业实现必须设填充比守卫，超限即拒绝。
\end{definition}

\begin{gapnote}
当前稀疏 Kron 采用最小 front（邻居数）贪心排序，属最小度思想的一个简化变体；
AMD 与嵌套分割排序尚未实现（手册"后续改进方向"已列为规划）。最小度族排序在
装配层 LU 求解器中的应用由求解后端（SuiteSparse/Eigen）自带，不在图模块内。
\end{gapnote}

\subsection{多级降阶的映射代数与等值电路}\label{sec:gr-th-compose}

\begin{proposition}[多级降阶的映射复合与逆序恢复]\label{sec:gr-th-compose-prop}
设降阶序列 $R_1,R_2,\dots,R_m$ 依次作用于网络，第 $k$ 步的状态空间由
$X^{(k-1)}$ 压缩到 $X^{(k)}$。若每步保留一个线性恢复算子
$T_k:X^{(k)}\to X^{(k-1)}$（无源消去时
$x^{(k-1)}=\begin{bmatrix}I\\-K_k\end{bmatrix}x^{(k)}$ 的嵌入形式；
收缩与重编号时为置换/复制算子），则全状态恢复为算子复合
\begin{equation}
x^{(0)}=T_1\,T_2\,\cdots\,T_m\,x^{(m)},
\end{equation}
即恢复必须按降阶的\emph{逆序}逐层施加。节点重编号对应置换同余
$Y\mapsto PYP^{\mathsf T}$，不改变任何谱与解的性质。
\end{proposition}
\begin{proof}
对 $m$ 归纳。$m=1$ 即单层恢复算子定义。设前 $m-1$ 层恢复为
$x^{(0)}=T_1\cdots T_{m-1}x^{(m-1)}$，而第 $m$ 层给出
$x^{(m-1)}=T_mx^{(m)}$，代入即得。每层算子由该层降阶类型唯一确定
（Kron 为回代算子，收缩为代表元复制，串联/悬垂为逐节点递推），复合良定义。
重编号是节点集的双射，$P$ 为正交置换阵，$PYP^{\mathsf T}$ 与原矩阵相似同构，
方程组仅换元不解变。
\end{proof}

\begin{proposition}[零阻抗等电位合并的精确性]\label{sec:gr-th-zmerge}
设边集合 $E_0$ 中每条边阻抗为零（闭合开关/断路器的理想化）。将每个由 $E_0$
连通的节点组并为一个超节点（电压取公共值、注入求和、外部支路平移端点），
则等值网络与原网络在同一边界条件下同解。
\end{proposition}
\begin{proof}
零阻抗边的 KVL 给出其两端电压相等，故组内所有节点电压必取同一值，以超节点
变量代替不损失自由度；组内节点的 KCL 方程求和，内部零阻抗边电流成对抵消，
得超节点注入等于组内注入之和；跨组支路端点平移后方程形式不变。反向，由超节点
解把公共电压复制回组内各节点，内部边电流不确定（零阻抗边潮流不可由电压恢复），
但全部节点量与原方程组相容。故对节点电压与外部端口行为，合并精确。
\end{proof}

\begin{theorem}[串联等值：度 2 无源节点的消去]\label{sec:gr-th-series}
设节点 $j$ 度为 2，经阻抗 $z_1$ 接 $i$、经 $z_2$ 接 $k$，$j$ 无注入、
无对地支路。消去 $j$ 等价于以单支路 $z_{eq}=z_1+z_2$ 直连 $i$--$k$，
即导纳
\begin{equation}
y_{eq}=\frac{y_1y_2}{y_1+y_2}.
\end{equation}
\end{theorem}
\begin{proof}
记串联电流 $i_1$（$i\to j$）、$i_2$（$j\to k$）。KVL：
$V_i-V_j=z_1i_1$，$V_j-V_k=z_2i_2$；$j$ 无注入无对地支路，KCL 给出
$i_1=i_2=:i$。两式相加消去 $V_j$：$V_i-V_k=(z_1+z_2)i$，恰为阻抗
$z_1+z_2$ 的单支路方程。此即为 $|\beta|=1$ 的 Kron 消去
（定理~\ref{sec:gr-th-kron-exact}）的显式闭式。
\end{proof}

\begin{remark}[串联等值的适用边界]
定理要求中间节点\emph{无注入且无对地元件}；若 $j$ 带负荷、发电机或线路充电
（$\pi$ 模型把充电放在节点上），串联合并必须迁移注入或改为近似处理。含
$\pi$ 充电的两段线路严格等值不是裸阻抗相加，而是含集中参数的 II 型等值；
实现层据此对串联候选施加设备类别与注入约束。
\end{remark}

\begin{proposition}[星-网变换即星形中心的 Kron 消去]\label{sec:gr-th-starmesh}
设星形中心节点 $0$ 经导纳 $y_\ell$ 连接叶节点 $\ell=1,\dots,n$，$0$ 无注入。
记 $y_\Sigma=\sum_{\ell=1}^n y_\ell$。消去 $0$ 后叶节点两两之间出现新支路
\begin{equation}
y_{\ell m}=\frac{y_\ell y_m}{y_\Sigma}\qquad(\ell\ne m),
\end{equation}
叶 $\ell$ 的对地自感纳增量使新对角元为 $y_\ell(y_\Sigma-y_\ell)/y_\Sigma$。
$n=3$ 时即经典星-三角（Y-$\Delta$）变换公式。
\end{proposition}
\begin{proof}
对 $|\beta|=1$ 应用 Schur 补：$Y_{00}=y_\Sigma$，$Y_{0\ell}=-y_\ell$，
\begin{equation}
[Y_{red}]_{\ell m}=Y_{\ell m}-\frac{Y_{\ell 0}Y_{0m}}{Y_{00}}
=0-\frac{y_\ell y_m}{y_\Sigma}\cdot(-1)^{2}
=-\frac{y_\ell y_m}{y_\Sigma}\quad(\ell\ne m),
\end{equation}
非对角元取负即新支路导纳 $y_{\ell m}=y_\ell y_m/y_\Sigma$；对角元
$y_\ell-y_\ell^2/y_\Sigma=y_\ell(y_\Sigma-y_\ell)/y_\Sigma$。
这正是消去图的成团规则（定理~\ref{sec:gr-th-fill}）：星形中心消去使其 $n$ 个
邻居成团，一次引入 $\binom{n}{2}$ 条新边——星-网变换是 fill-in 现象的最小
完整实例。
\end{proof}

\begin{remark}[悬垂折叠的近似性与工作点依赖]
度 1 叶节点 $\ell$ 经阻抗 $z_\ell$ 挂在父节点 $p$ 上。把叶负荷 $S_\ell$ 折算到
$p$ 的严格式应计入支路损耗
\begin{equation}
\Delta S_{\mathrm{loss}}=z_\ell\,\frac{|S_\ell|^2}{|V_\ell|^2},
\end{equation}
而 $|V_\ell|$ 又依赖全网络工况。折叠时若取 $|V_p|\approx1.0$ pu 估计损耗
（或忽略损耗），所得等值负荷的误差正比于 $|S_\ell|^2$——轻载时微不足道，
重载长馈线时平方放大。这类近似等值的误差\emph{内生地依赖工作点}，上报结果时
应声明近似口径而非宣称精确。
\end{remark}

\subsection{与实现的对应}
\begin{center}\footnotesize
\begin{longtable}{@{}L{6.8cm}L{8.6cm}@{}}
\toprule
理论条目 & 实现状态\\
\midrule
\endfirsthead
\toprule
理论条目 & 实现状态\\
\midrule
\endhead
图容器、邻接表与 domain-qualified 映射 &
\implfull{src/graph/power\_system\_graph.cpp:build\_power\_system\_graph}（结构定义见 include/hacdcpf/graph/power\_system\_graph.hpp）\\
关联矩阵分解 $Y_{bus}=AY_bA^{\mathsf T}$ 与正定性 &
\implpart{图模块只建拓扑不装配 $Y_{bus}$；装配在装配层完成（见《潮流计算技术手册》），本理论为其依据}\\
连通分量枚举与孤岛判定 &
\implfull{src/graph/topology\_analysis.cpp:analyze\_topology}（快速口径见同文件 count\_islands、is\_connected）\\
辐射性判据（连通 $\land$ $E=N-1$，推论~\ref{sec:gr-th-radial}） &
\implfull{src/graph/topology\_analysis.cpp:is\_radial}（$\fld{cycle\_count}=\max(0,E-N+p)$ 在 analyze\_topology）\\
基本回路基（定理~\ref{sec:gr-th-cyclebasis}） &
\implfull{src/graph/topology\_analysis.cpp:find\_fundamental\_cycles}（BFS 生成树 + 父指针回溯至 LCA）\\
Tarjan 桥/割点（引理~\ref{sec:gr-th-bridge}、\ref{sec:gr-th-ap}） &
\implfull{src/graph/topology\_analysis.cpp:tarjan\_bridge\_ap}（桥判据 $low[v]>disc[parent]$、非根割点 $low[v]\ge disc[parent]$ 与定理一致）\\
零阻抗等电位合并（命题~\ref{sec:gr-th-zmerge}）与并查集 &
\implfull{src/graph/switch\_contraction.cpp:contract\_zero\_impedance\_edges}（内部 DSU，路径压缩 + 按秩合并；含电压基与多 slack 守卫）\\
Kron 精确降阶（$I_\beta=0$，定理~\ref{sec:gr-th-kron-exact}） &
\implfull{src/graph/kron\_reduction.cpp:apply\_kron\_reduction}（FullPivLU::isInvertible 主元守卫 + 填充比守卫）\\
Ward 修正注入 $I_{red}$（定理~\ref{sec:gr-th-ward}） &
\implfull{src/graph/kron\_reduction.cpp:apply\_kron\_reduction\_with\_injection}\\
含 $I_\beta$ 的完整电压恢复 $V_\beta=-KV_\alpha+Y_{\beta\beta}^{-1}I_\beta$ &
\implpart{recover\_eliminated\_voltages 仅返回 $-KV_\alpha$，无 $I_\beta$ 修正项；调用链 result\_recovery.cpp:recover\_kron\_eliminated\_buses 同受此限}\\
消去块非奇异的充分条件（引理~\ref{sec:gr-th-ybb}） &
\implpart{理论充分条件未做静态判定；以运行时 FullPivLU 可逆性检查代替（apply\_kron\_reduction）}\\
稀疏逐节点 Kron 与成团更新（定理~\ref{sec:gr-th-fill} 的更新式） &
\implfull{src/graph/sparse\_kron\_reduction.cpp:reduce\_sparse\_kron}（最小 front 贪心 + front/主元/填充三重守卫，eliminate 保存恢复系数）\\
最小度/AMD/嵌套分割排序 &
\implpart{稀疏 Kron 内部为最小 front 简化变体；AMD/嵌套分割未实现（"后续改进方向"已列）；装配层 LU 排序由求解后端自带}\\
最小填入的 NP 完全性结论 &
\implnone（纯理论结论，无需实现）\\
多级映射复合与逆序恢复（命题~\ref{sec:gr-th-compose-prop}） &
\implfull{src/graph/result\_recovery.cpp:recover\_kron\_eliminated\_buses}（映射结构 ReductionMapping 见 include/hacdcpf/graph/reduction\_mapping.hpp；稀疏恢复算子见 sparse\_kron\_recovery\_operator）\\
串联等值 $z_1+z_2$（定理~\ref{sec:gr-th-series}） &
\implpart{apply\_series\_reduction 受设备类别、保留集、注入与映射可恢复性约束，非裸公式逐式实现（见源码等价章）}\\
星-网/Y-$\Delta$ 变换（命题~\ref{sec:gr-th-starmesh}） &
\implpart{未单列变换例程；其机理已由稀疏 Kron 的成团消去涵盖}\\
悬垂折叠近似与工作点依赖 &
\implpart{apply\_pendant\_reduction 为近似（$|V_{\text{parent}}|\approx1.0$ pu）；带实际父电压的迭代恢复在 result\_recovery.cpp:recover\_pendant\_buses，恢复收敛标志未上报}\\
REI/Dimo 等高级外部等值 &
\implnone\\
\bottomrule
\end{longtable}
\end{center}

\subsection{参考文献}
{\renewcommand{\section}[2]{}%
\begin{thebibliography}{9}\small
\bibitem{grf:tarjan} R.~E. Tarjan, Depth-first search and linear graph algorithms,
\emph{SIAM Journal on Computing}, 1972.
\bibitem{grf:rose} D.~J. Rose, Triangulated graphs and the elimination process,
\emph{Journal of Mathematical Analysis and Applications}, 1970.
\bibitem{grf:tinney} W.~F. Tinney and J.~W. Walker, Direct solutions of sparse
network equations by optimally ordered triangular factorization,
\emph{Proceedings of the IEEE}, 1967.
\bibitem{grf:yannakakis} M.~Yannakakis, Computing the minimum fill-in is
NP-complete, \emph{SIAM Journal on Algebraic and Discrete Methods}, 1981.
\bibitem{grf:georgeliu} A.~George and J.~W.-H. Liu, \emph{Computer Solution of
Large Sparse Positive Definite Systems}, Prentice-Hall, 1981.
\bibitem{grf:kron} G.~Kron, \emph{Diakoptics: The Piecewise Solution of
Large-Scale Systems}, MacDonald, 1963.
\bibitem{grf:ward} J.~B. Ward, Equivalent circuits for power-flow studies,
\emph{Transactions of the AIEE}, 1949.
\bibitem{grf:arrillaga} J.~Arrillaga and N.~R. Watson, \emph{Computer Modelling
of Electrical Power Systems}, 2nd ed., Wiley, 2001.
\bibitem{grf:even} S.~Even, \emph{Graph Algorithms}, 2nd ed., Cambridge
University Press, 2012.
\bibitem{grf:zhang} 张伯明、陈寿孙、严正，《高等电力网络分析》（第二版），
清华大学出版社，2007.
\end{thebibliography}}

% ── 实现转写层（源码等价）──
\section{源码等价图算法}

本章只描述当前图实现实际计算的量，不用理论上更完整的降阶式替换代码接口。

\subsection{拓扑量}

\srcpath{src/graph/topology\_analysis.cpp:count\_islands} 只统计 in-service 节点和边。
空节点图的 \srcpath{is\_radial} 返回真；非空图只有在 \srcpath{count\_islands()==1} 且
$E=N-1$ 时返回真。完整报告先以 DFS 得到分量数 $p$，再赋值
\begin{equation}
 \fld{cycle\_count}=\max(0,E-N+p),\qquad
 \fld{is\_radial}=(p=1)\land(\fld{cycle\_count}=0).
\end{equation}
AC/DC 混合岛只要包含任何 AC 节点就归为 AC 岛；AC 岛要求 AC slack，纯 DC 岛要求 DC 电压参考。
桥判据为 DFS 子节点 $v$ 满足 $low[v]>disc[parent]$；非根割点为
$low[v]\ge disc[parent]$，根割点要求 DFS 子树数大于 1。
依据为 \srcpath{src/graph/topology\_analysis.cpp:tarjan\_bridge\_ap} 和
\srcpath{src/graph/topology\_analysis.cpp:analyze\_topology}。

\subsection{稠密 Kron 降阶}

\srcpath{src/graph/kron\_reduction.cpp:apply\_kron\_reduction} 把输入稀疏矩阵累加到四个稠密块。
消去集为空时仅抽取 retained 子矩阵并直接标记 valid。
否则使用 \srcpath{FullPivLU::isInvertible()} 拒绝奇异 $Y_{\beta\beta}$，再计算
\begin{align}
 K&=Y_{\beta\beta}^{-1}Y_{\beta\alpha},\\
 Y_{red}&=Y_{\alpha\alpha}-Y_{\alpha\beta}K.
\end{align}
代码还显式计算并保存 $Y_{\alpha\beta}Y_{\beta\beta}^{-1}$。
非零数用 $|x|>10^{-15}$ 计数，填充比为
\begin{equation}
 r_{fill}=\begin{cases}
 nnz(Y_{red})/nnz(Y_{full}),&nnz(Y_{full})>0,\\1,&nnz(Y_{full})=0.
 \end{cases}
\end{equation}
只有 $r_{fill}\le\fld{max\_fill\_ratio}$ 才输出 reduced 矩阵并标记 valid。

带注入入口在上述 valid 后计算
\begin{equation}
 I_{red}=I_\alpha-\left(Y_{\alpha\beta}Y_{\beta\beta}^{-1}\right)I_\beta,
\end{equation}
见 \srcpath{src/graph/kron\_reduction.cpp:apply\_kron\_reduction\_with\_injection}。但是当前恢复 API
\srcpath{recover\_eliminated\_voltages} 不接收 $I_\beta$，实际只返回
\begin{equation}
 V_\beta=-K V_\alpha.
\end{equation}
invalid 或 $K$ 行数为零时返回空向量，见 \srcpath{src/graph/kron\_reduction.cpp:recover\_eliminated\_voltages}。
因此不能声称当前 API 已实现含非零 $I_\beta$ 的完整电压恢复式。

\subsection{稀疏逐节点 Kron}

\srcpath{src/graph/sparse\_kron\_reduction.cpp:eliminate} 对节点 $k$ 的每个列邻居 $i$、行邻居 $j$
执行
\begin{equation}
 Y_{ij}\leftarrow Y_{ij}-\frac{Y_{ik}Y_{kj}}{Y_{kk}},
\end{equation}
并保存恢复系数 $-Y_{kj}/Y_{kk}$。任何更新值绝对值不大于
\fld{drop\_tolerance} 就从 row/column map 中删除。

候选节点按 $\max(n_{row},n_{col})$ 的 front size 最小优先；弹出时若 front 已变化则重新入队。
只有 front 不超过 \fld{max\_front}、$|Y_{kk}|$ 严格大于 \fld{pivot\_tolerance}，且预测消去后
非零数不超过
\begin{equation}
 \left\lceil\fld{max\_nnz\_ratio}\cdot nnz(Y_{initial})\right\rceil
\end{equation}
才消去。选项要求 max front 非负、max nnz ratio 至少 1、两个 tolerance 非负；否则返回 error。
依据为 \srcpath{src/graph/sparse\_kron\_reduction.cpp:reduce\_sparse\_kron}。

\subsection{串联、悬垂与映射覆盖边界}

串联和悬垂降阶不是单一 $Z_1+Z_2$ 公式：代码还受设备类别、保留集合、注入、状态和映射
可恢复性约束。其当前实现分别位于 \srcpath{src/graph/series\_reduction.cpp}、
\srcpath{src/graph/reduction\_plan.cpp} 与 \srcpath{src/graph/result\_recovery.cpp}。
本章未逐分支展开这些文件，因此不提供“等值阻抗相加”或“叶节点固定 1 pu”等概括式，也不声称
这些路径已完成数学逐式覆盖。


\section{算法流程}
\begin{figure}[H]\centering
\begin{tikzpicture}[font=\footnotesize,
  box/.style={draw,rounded corners,minimum height=0.85cm,align=center,inner sep=3pt},arr/.style={->,thick}]
  \node[box](build) at (0,0){建图\\build\_power\_system\_graph};
  \node[box](topo) at (3.3,0){拓扑分析\\岛/桥/割点/回路};
  \node[box](plan) at (6.6,0){降阶计划\\make\_reduction\_plan};
  \node[box](apply) at (9.9,0){逐级降阶\\收缩/串联/悬垂/Kron};
  \node[box](solve) at (3.3,-1.9){装配求解\\$Y_{\text{red}}$};
  \node[box](rec) at (6.6,-1.9){状态恢复\\recover\_*};
  \node[box](full) at (9.9,-1.9){FullNetworkVoltages\\AC/DC 分列};
  \draw[arr](build)--(topo);\draw[arr](topo)--(plan);\draw[arr](plan)--(apply);
  \draw[arr](apply.south)|-(solve.east);\draw[arr](solve)--(rec);\draw[arr](rec)--(full);
\end{tikzpicture}
\caption{图建模—降阶—求解—恢复流水线}
\end{figure}
降阶计划 \srcpath{make\_reduction\_plan} 先分类候选（可收缩边、可串联节点、悬垂叶、
Kron 消去集），逐类应用并在每步追加映射；求解后按逆序恢复。桥/割点由单趟 Tarjan DFS
在 $O(N+E)$ 内求出。

\section{工程假设与数值稳定性分析}
\subsection{工程假设}
无向图假设支路双向对称；串联合并假设中间节点无横向注入；Kron 假设消去节点为无源或
线性；悬垂折叠假设父节点电压近似 $1.0$ pu。以上假设违反时相应降阶被守卫拒绝或标注为近似。

\subsection{数值稳定性}
Kron 的数值风险集中在 $Y_{\beta\beta}^{-1}$：接近奇异（弱耦合子网）会放大恢复误差，故设
主元守卫；填充比守卫防止稀疏 Kron 退化为稠密（内存与时间爆炸）。开关收缩为纯拓扑操作
（无浮点求逆），数值稳定；串联合并只做复数加法，稳定。恢复映射为复合而非覆盖，避免
多级降阶时的编号错位。

\section{复杂度分析}
建图 $O(N+E)$；拓扑分析（连通分量、桥、割点、基本圈）单趟 DFS $O(N+E)$；开关收缩为
并查集 $O(E\,\alpha(N))$；串联/悬垂降阶按候选数线性；稠密 Kron 为 $O(|\beta|^3)$
（$Y_{\beta\beta}$ 求逆），稀疏 Kron 在良好排序下接近 $O(\text{nnz})$ 但受填充比限制；
状态恢复为一次前代/回代，与降阶同阶。

\section{异常工况处理}
\begin{itemize}
  \item \textbf{$Y_{\beta\beta}$ 奇异}：主元守卫触发，拒绝该 Kron 降阶并回退到未降阶网络；
  \item \textbf{填充比超限}：\fld{max\_fill\_ratio} 守卫拒绝降阶，保持稀疏；
  \item \textbf{多平衡节点欲合并}：默认拒绝并诊断，除非显式 \fld{allow\_multi\_slack\_merge}；
  \item \textbf{电压基不一致的零阻抗合并}：拒绝合并并告警（避免跨电压等级等电位错误）；
  \item \textbf{死岛（无 slack 的孤岛）}：拓扑报告标注，上层据此置零或跳过该岛。
\end{itemize}

\section{与其他模块的接口关系}
图模块为投影层（\srcpath{src/model/network\_utils.cpp} 的零阻抗合并）与装配层
（$Y_{bus}$ 构建）提供拓扑基础；潮流/OPF 在装配前调用降阶、求解后调用恢复；重构
（\srcpath{network\_reconfiguration}）与弹性（\srcpath{resilience}）消费辐射性/桥/割点判定；
可靠性 N-1 消费割点与桥。domain-qualified 映射是全平台跨 AC/DC 传递的统一契约。

\section{测试与验证建议}
注册测试：\srcpath{test\_graph}（拓扑、开关收缩、串联/悬垂）、
\srcpath{test\_graph\_kron}（Schur 补与电压恢复）、
\srcpath{test\_graph\_roundtrip}（聚合前后 PF/OPF 一致性 round-trip）。建议补充：
AC/DC 同号 ID 的 domain-qualified 恢复对照、多级降阶复合映射逆运算、$Y_{\beta\beta}$
近奇异时守卫触发、填充比守卫边界、以及死岛与多岛混合的拓扑断言。

\section{工业级评估指标}
建议纳入：降阶后 $Y_{bus}$ 维数压缩率与求解加速比；Kron 恢复的电压相量最大误差
（对无源节点应达机器精度量级）；开关收缩后端子潮流恢复的守恒残差；拓扑分析对已知
桥/割点算例的召回率；round-trip（降阶—求解—恢复）与直接求解的逐母线一致性（$L_\infty$ 范数）。

\section{适用范围与局限性}
\begin{itemize}
  \item Kron 仅对无源/线性内部节点精确；恒功率内部节点为工作点相关的近似；
        \fld{max\_fill\_ratio} 填充守卫与奇异 $Y_{\beta\beta}$ 主元守卫可拒绝降阶；
  \item 悬垂折叠为近似（取 $|V_{\text{parent}}|=1.0$ pu；迭代恢复不单独报告收敛标志）；
  \item 旧式扁平映射（\fld{bus\_id\_to\_node\_idx}、\texttt{IslandInfo::bus\_ids}、
        \texttt{ReductionMapping} 扁平映射）在 AC/DC 同号 ID 时含糊，须用 domain-qualified 映射；
  \item 多平衡节点合并默认拒绝，除非 \fld{allow\_multi\_slack\_merge}；两端电压基须在容差内一致。
\end{itemize}

\section{后续改进方向}
\begin{enumerate}
  \item 为恒功率内部节点的 Kron 近似提供工作点相关的误差界估计并写入结果；
  \item 悬垂折叠的迭代恢复已实现（带父节点实际电压，\srcpath{result\_recovery.cpp}），
        剩余工作是把恢复收敛标志写入结果并上报；
  \item 引入基于最小度/近似最小填充（AMD）的消去排序以进一步压低稀疏 Kron 填充；
  \item 统一淘汰旧式扁平映射，全平台强制 domain-qualified 接口。
\end{enumerate}

\section{跨模块工业级技术评价基线}

本章规定本手册所述模块进入工程算例、批量研究或生产服务前必须共同满足的评价基线。
具体模型公式、变量、约束和专用算法以正文相应章节为准；本章不扩大源码已经实现的范围。

\subsection{输入、输出、边界条件与数据结构审计}
输入审计应逐项检查有限性、单位、上下界、枚举值和引用完整性。系统对象必须先按所需层级完成
校验；AC 与 DC 母线采用域限定映射，组件稳定 \texttt{index}、向量位置和临时图索引不得混用。
输出审计必须记录单位、索引空间、模型范围、收敛或可行状态、后端、容差、迭代/节点计数和告警。
空系统、空时域、孤岛、重复 ID、零容量、退化约束与不可用可选后端均属于边界测试集。

\subsection{工程假设与数值稳定性分析}
模型使用的平衡、线性化、稳态、独立性、概率分布、时间离散或网络等值假设必须与应用场景匹配。
数值评价至少包含变量缩放、绝对/相对残差、可行性违反、条件数或枢轴质量、步长/阻尼、迭代停滞
和随机种子复现性。若采用正则化、平滑、回退、松弛或近似，应同时报告触发条件、参数值、对主指标
的敏感性以及是否改变模型口径；不得用“已返回结果”替代收敛或有效性证明。

\subsection{复杂度与资源分析}
令 $n_x$、$n_c$、$n_t$、$n_s$、$|V|$、$|E|$ 分别表示变量、约束、时段、场景、节点和边数。
报告应分别给出建模、求解、后处理的时间与内存。图遍历通常为 $O(|V|+|E|)$；逐时段/逐场景
流程至少线性增长为 $O(n_t n_s)$ 次子问题调用；稠密线性代数最坏可达 $O(n_x^3)$，稀疏方法则应
报告结构非零数、填充率和因子分解峰值内存；MILP/NLP 不给出虚假的多项式最坏界，而报告节点数、
间隙、时限和实测扩展曲线。复杂度结论必须基于固定硬件、固定后端和可复现算例族。

\subsection{异常工况处理}
非法输入应在进入求解器前返回结构化错误；不可行、无界、不收敛、时限、内存不足和后端缺失必须
相互区分。部分结果只有在对应 \fld{ValidityFlags}、\fld{model\_scope}、
\fld{model\_limitations} 或等价字段允许时才可使用。异常路径不得静默丢弃 AC/DC 资产、制造平衡
节点、把 NaN 改写为零或把近似解标为全局最优；系统原始对象应保持可恢复和可重试。

\subsection{与其他模块的接口关系}
接口链路遵循“工程语义模型 $\to$ 校验 $\to$ 投影/装配 $\to$ 求解或分析 $\to$ 结果回投影”。
跨模块传递时保留稳定 ID、单位、符号、时间基准和场景身份；消费潮流、OPF、动态或可靠性结果前，
先检查其收敛、覆盖范围和有效性标志。HTTP/JSON/Python 层只做语义保持适配，不重新解释求解结果。

\subsection{测试与验证建议}
最低验证矩阵包括：手算小例、标准算例、边界/异常输入、随机性质测试、算法或后端交叉校核、
序列化往返、稳定 ID 回投影、AC/DC 同号母线、重复运行确定性和规模扩展测试。数值算法还应加入
残差独立重算、雅可比有限差分或自动微分校核；近似模型应与高保真模型比较并给出误差分布，而非
只报告单个平均值。回归报告必须列明构建配置、算例版本、容差、通过/失败/跳过数及未验证范围。

\subsection{工业级评估指标}
\begin{itemize}
  \item \textbf{正确性}：守恒残差、约束最大违反、解析/基准误差、结果回投影一致性；
  \item \textbf{鲁棒性}：收敛率、不可行识别率、异常分类准确率、扰动敏感性与随机复现率；
  \item \textbf{性能}：端到端时延、吞吐量、峰值内存、并行效率与规模增长斜率；
  \item \textbf{可追溯性}：输入模式版本、稳定 ID 覆盖率、模型范围/限制字段完整率、告警可定位率；
  \item \textbf{工程有效性}：与标准、外部引擎或现场数据的偏差，以及误判/漏判对业务指标的影响。
\end{itemize}
指标应预先给出验收阈值和失败处置；未达到阈值时只能标记为实验性或受限适用。

\subsection{适用范围、局限性与后续改进}
本手册只评价当前源码覆盖的模型、后端和数据结构，不对未建模物理过程、未安装求解器或未执行的
外部验证作保证。后续改进应按“缺口 $\to$ 可量化指标 $\to$ 源码/测试变更 $\to$ 独立验证”闭环，
优先消除会造成错误结论的语义缺口，其次提升数值鲁棒性与性能，最后扩展模型范围。


\end{document}
